\BourbakiExerciseGroup

\begin{enumerate}
\def\labelenumi{\arabic{enumi})}
\tightlist
\item
设 \((x_{t})\) 是一个实数族，其所有项都属于区间 \(A' = ] - \infty, +\infty]\) （相应地 \(A'' = [-\infty, +\infty[\) ）。族 \((x_{t})\) 在 \(\overline{\mathbf{R}}\) 中可和，当且仅当下列条件之一成立：
\end{enumerate}

\begin{enumerate}
\def\labelenumi{\alph{enumi})}
\item
数 \(x_{t}\) 中至少有一个等于 \(+\infty\) （相应地 \(-\infty\) ）。
\item
两个和 \(\sum_{i} x_{i}^{+}, \sum_{i} x_{i}^{-}\) 中至少有一个是有限的。
\end{enumerate}

在第一种情形有 \(\sum_{i} x_{i} = +\infty\) （相应地 \(-\infty\) ）；在第二种情形有 \(\sum_{i} x_{i} = \sum_{i} x_{i}^{+} - \sum_{i} x_{i}^{-}\) 。

\begin{enumerate}
\def\labelenumi{\arabic{enumi})}
\setcounter{enumi}{1}
\item
设 \((x_{i})\) 是一个实数族，其所有项都属于区间 \(\mathbf{A}'\) （相应地 \(\mathbf{A}''\) ），且满足习题 1 的条件 \(b\) )。试证：\((x_{i})\) 的每个子族在 \(\overline{\mathbf{R}}\) 中可和，并且 \(x'\) 的和是可结合的（参看第三章，§ 5，第 3 节，定理 2）。
\item
设 \((x_{n})_{n\geqslant 0}\) 是有限实数 \(\geqslant 0\) 的一个递减序列。这个序列在 \(\mathbf{R}\) 中可和，当且仅当序列 \((2^{n}x_{2^{n}})_{n\geqslant 0}\) 在 \(\mathbf{R}\) 中可和（``Cauchy 并项判别法''）。* 由此推出：通项为 \(1 / n^{\alpha}\) 的级数当 \(\alpha >1\) 时收敛，当 \(0 < \alpha \leqslant 1\) 时不收敛。*
\item
设 \((a_{n})\) 是有限实数 \(\geqslant 0\) 的任一序列。试证：序列 \(\left(\frac{a_n}{(1 + a_1)(1 + a_2)\ldots(1 + a_n)}\right)\) 在 \(\mathbf{R}\) 中可和（把这个序列的每一项表示为一个差）。
\item
设 \((d_n)\) 是有限实数 \(\geqslant 0\) 的一个序列，满足
\end{enumerate}

\[
\sum_ {n = 0} ^ {\infty} d _ {n} = + \infty .
\]

关于通项分别为下列各量的级数的收敛性，能说些什么？

\[
\frac {d _ {n}}{\mathrm{I} + d _ {n}}, \quad \frac {d _ {n}}{\mathrm{I} + n d _ {n}}, \quad \frac {d _ {n}}{\mathrm{I} + n ^ {2} d _ {n}}, \quad \frac {d _ {n}}{\mathrm{I} + d _ {n}}?
\]

\begin{enumerate}
\def\labelenumi{\arabic{enumi})}
\setcounter{enumi}{5}
\item
试证：\(s = \sum_{n=1}^{\infty} (1/n)\) 等于 \(+\infty\) ，办法是证明 \(s \geqslant s + \frac{1}{2}\) （求出偶指标各项之和与奇指标各项之和的、以 \(s\) 为函数的下界）。
\item
试证：对每个整数 \(p > \mathbf{r}\) ，有
\end{enumerate}

\[
\sum_ {n = 1} ^ {\infty} \frac {(- \mathrm{I}) ^ {n - 1}}{n ^ {p}} = (\mathrm{I} - 2 ^ {1 - p}) \sum_ {n = 1} ^ {\infty} \frac {\mathrm{I}}{n ^ {p}}.
\]

\begin{enumerate}
\def\labelenumi{\arabic{enumi})}
\setcounter{enumi}{7}
\item
试证：若 \(m\) 是一个整数 \(> 0\) ，则 \(\sum_{n \geqslant 1, n \neq m} \frac{1}{m^2 - n^2} = -\frac{3}{4m^2}\) （把有理分式 \(\frac{1}{m^2 - x^2}\) 表示为部分分式之和）。
\item
若通项为 \(x_{n}\) 的级数在 \(\mathbf{R}\) 中收敛，试证：\(\liminf_{n\to \infty}nx_n\leqslant 0\leqslant \limsup_{n\to \infty}nx_n\) 。
\end{enumerate}

¶ 10) 级数 \((u_n)\) 在 \(\mathbb{R}\) 中收敛，当且仅当：对每个递增序列 \((p_n)\) （数 \(>0\) ，满足 \(\lim_{n\to \infty}p_n = +\infty\) ），有 \(\lim_{n\to \infty}\left(\left(\sum_{k=0}^{n}p_ku_k\right) / p_n\right) = 0\) （用 § 5 习题 15）。

\begin{enumerate}
\def\labelenumi{\arabic{enumi})}
\setcounter{enumi}{10}
\tightlist
\item
考察一个序列 \((a_{n})\) ，它的每一项都是有限实数的一个有限序列之和
\end{enumerate}

\[
a _ {n} = b _ {n, 1} + b _ {n, 2} + \dots + b _ {n, k _ {n}}.
\]

对每对正整数 \((n, p)\) （满足 \(p \leqslant k_n\) ），若 \(m = \sum_{i=0}^{\infty} k_i + p\) ，令 \(c_m = b_{n,p}\) 。试证：若通项为 \(a_n\) 的级数在 \(\mathbf{R}\) 中收敛，且

\[
d _ {n} = \left| b _ {n, 1} \right| + \left| b _ {n, 2} \right| + \dots + \left| b _ {n, k _ {n}} \right|
\]

当 n 趋于无穷时趋于 o，则通项为 \(c_{m}\) 的级数收敛，且与通项为 \(a_{n}\) 的级数有相同的和。

¶ 12) 设 \((u_n)\) 是有限实数的一个序列。对集合 N 的每个置换 \(\sigma\) ，令

\[
r (n) = | \sigma (n) - n |. \sup _ {m \geqslant n} | u _ {m} |.
\]

\begin{enumerate}
\def\labelenumi{\alph{enumi})}
\item
若通项为 \(u_n\) 的级数收敛，且 \(\lim_{n \to \infty} r(n) = 0\) ，试证：通项为 \(v_n = u_{\sigma(n)}\) 的级数收敛到同一个和。{[}考察差 \(\sum_{k=0}^{n} v_k - \sum_{k=0}^{n} u_k\) （对大的 \(n\) 值），且若 \(h\) 是最小的整数 \(\leqslant n\) （它满足 \(\sigma(h) > n\) ），注意在这个差中，两个和的每一个里至多有 \(\sigma(h) - h\) 项不被消去{]}。
\item
设通项为 \(u_n\) 的级数收敛。试给出一个置换 \(\sigma\) 的例子，使得 \(\lim_{n \to \infty} r(n) = 0\) ，但不满足第三章，§ 5，习题 6 a) 的判别法。{[}定义一族两两不相交的区间 \(I_k = [n_k - 2k, n_k + 2k]\) （在 \(N\) 中），并取 \(\sigma\) 使得 \(\sigma(n) = n\) 每当 \(n\) 不属于任何 \(I_k\) 时成立，且 \(\sigma(n_k - 2j) = n_k + 2j\) ， \(\sigma(n_k + 2j) = n_k - 2j\) 对一切 \(k\) 与 \(0 \leqslant j \leqslant k\) 成立。{]}
\end{enumerate}

\begin{enumerate}
\def\labelenumi{\arabic{enumi})}
\setcounter{enumi}{12}
\item
设 \((u_n)\) 是实数 \(\geqslant 0\) 的一个序列，满足 \(\lim_{n\to \infty}u_n = 0\) 且 \(\sum_{n = 0}^{\infty}u_n = +\infty\) 。设 \((p_n)\) 是数 \(>0\) 的一个严格递增序列，满足 \(\lim_{n\to \infty}p_n = +\infty\) 。试证：存在一个置换 \(\sigma\) （ \(\mathbf{N}\) 的），使得对每个 \(n\in \mathbf{N}\) ， \(\sum_{k = 0}^{n}u_{\sigma (n)}\leqslant p_n\)
\item
设 \((u_{n})\) 是有限实数的一个序列，使得通项为 \(u_{n}\) 的级数收敛但不绝对收敛；令 \(s = \sum_{n=0}^{\infty} u_{n}\) 。试证：对每个数 \(s' \geqslant s\) ，存在一个置换 \(\sigma\) （ \(\mathbf{N}\) 的），使得 \(\sigma(n) = n\) 对那些指标 \(n\) （满足 \(u_{n} \geqslant 0\) ）成立，并且 \(\sum_{n=0}^{\infty} u_{\sigma(n)} = s'\) 。{[}对 \(m\) 用归纳法证明：存在一个置换 \(\sigma_m\) （ \(\mathbf{N}\) 的），使得 \(\sigma_m(k) = k\) 对一切 \(k\) （满足 \(u_k \geqslant 0\) ）成立，并且若用 \(u_k^{(m)}\) 表示 \(u_{\sigma_m(k)}\) ，则有一个整数 \(p_m\) 具有如下性质：\(\left| s' - \sum_{i=0}^{k} u_i^{(m)} \right| \leqslant 1/m\) 对一切 \(k \geqslant p_m\) 成立。此外，\(\sigma_{m+1}\) 满足：\(\sigma_{m+1}(k) = \sigma_m(k)\) 对一切 \(k\) （满足 \(\sigma_m(k) < p_m\) ）以及一切 \(k\) （满足 \(u_k \leqslant -1/m\) ）成立。{]}
\item
设 \((u_{n})\) 是有限实数的一个序列，满足 \(\lim_{n\to \infty}u_n = 0\) 且 \(\sum_{n}u_{n}^{+} = \sum_{n}u_{n}^{-} = +\infty\) 。若 \(a\) 与 \(b\) 是任意两个实数（有限或无穷），满足 \(a\leqslant b\) ，试证：存在一个置换 \(\sigma\) （ \(\mathbf{N}\) 的），使得当令 \(s_n = \sum_{k = 0}^n u_{\sigma (k)}\) 时，有 \(\liminf_{n\to \infty}s_n = a\) 与 \(\limsup_{n\to \infty}s_n = b\) 。试证：序列 \((s_n)\) 的聚点所成之集于是就是区间 \([a,b]\) 。
\item
设 \((u_n)\) 是有限实数的一个序列，使得通项为 \(u_n\) 的级数收敛但不绝对收敛。试证：存在一个置换 \(\sigma\) （ \(\mathbf{N}\) 的），它满足第三章，§ 5，习题 6 \(a\) ) 的条件，但使通项为 \(u_{\sigma - 1(n)}\) 的级数不收敛。{[}设 \(h, k, m\) 是具有下列性质的三个整数：若 \(p_1, \ldots, p_r\) 是整数 \(n\) （满足 \(h \leqslant n < h + m\) 且 \(u_n \geqslant 0\) ），则 \(u_{p_1} + \cdots + u_{p_r} > 2\) 且 \(h + m < k - r\) ；此外若 \(\mu, \nu\) 是使 \(h + m \leqslant \mu \leqslant \nu\) 的任意整数，则 \(\left|\sum_{n = \mu}^{\nu} u_n\right| \leqslant 1\) 。令 \(s = m - r\) ，并用 \(q_1, \ldots, q_s\) 表示这些整数 \(n\) （它们满足 \(h \leqslant n < h + m\) 且 \(u_n < 0\) ）。考察一个置换 \(\pi\) ——区间 \([h, k + s]\) （ \(\mathbf{N}\) 中的）的置换——它满足：\(\pi(p_i) = k - i + 1\) （对 \(i \leqslant i \leqslant r\) ）， \(\pi(q_j) = k + j\) （对 \(i \leqslant j \leqslant s\) ）， \(\pi(k + j) = h + s - j\) （对 \(i \leqslant j \leqslant s\) ）， \(\pi(k - i + 1) = h + s + i - 1\) （对 \(i \leqslant i \leqslant r\) ），最后 \(\pi(n) = n\) （对 \(h + m \leqslant n \leqslant k - r\) ）；试证我们有
\end{enumerate}

\[
\left| \sum_ {i = h} ^ {k} u _ {i} - \sum_ {i = h} ^ {k} u _ {\pi^ {- 1} (i)} \right| \geqslant I.
\]

\begin{enumerate}
\def\labelenumi{\arabic{enumi})}
\setcounter{enumi}{16}
\tightlist
\item
若 \(u_{mn} = \mathbf{i} / (m^2 - n^2)\) （对 \(m \neq n\) ）且 \(u_{mm} = 0\) ，试证
\end{enumerate}

\[
\sum_ {m = 1} ^ {\infty} \left(\sum_ {n = 1} ^ {\infty} u _ {m n}\right) = - \sum_ {n = 1} ^ {\infty} \left(\sum_ {m = 1} ^ {\infty} u _ {m n}\right) \neq 0.
\]

（用习题 8。）

\begin{enumerate}
\def\labelenumi{\arabic{enumi})}
\setcounter{enumi}{17}
\item
设 \((\mathbf{i} + u_{\mathfrak{i}})\) 是一个实数族，其所有项都属于区间 \([0, +\infty[\) （相应地 \(]0, +\infty]\) ）。族 \((\mathbf{i} + u_{\mathfrak{i}})\) 在 \(\mathbb{R}\) 中可乘，当且仅当下列条件之一成立：a) 数 \(\mathbf{i} + u_{\mathfrak{i}}\) 中至少有一个等于 o（相应地 \(+\infty\) ）。b) 族 \(u_{\mathfrak{i}}\) 在 \(\overline{\mathbb{R}}\) 中可和。
\item
设 \((\mathbf{I} + u_{t})_{t\in \mathbf{I}}\) 是 \(\mathbb{R}\) 中的一个可乘族。对每个非空子集 \(\mathbf{H}\) （ \(\mathbf{I}\) 的），用 \(v_{\mathbf{H}}\) 表示 \(\prod_{i\in \mathbf{H}}u_i\) 。试证：族 \((v_{\mathbf{H}})_{\mathbf{H}\in \mathfrak{F}(\mathbf{I})}\) 在 \(\mathbb{R}\) 中可和，并且
\end{enumerate}

\[
\mathrm{I} + \sum_ {\mathbf {H}} v _ {\mathbf {H}} = \prod_ {i \in \mathbf {I}} (\mathrm{I} + u _ {i}).
\]

逆命题是否成立？

由此推出：若 \(-\mathbf{I} < x < \mathbf{I}\) ，则 \(\prod_{n=0}^{\infty} (\mathbf{I} + x^{2^n}) = \mathbf{I} / (\mathbf{I} - x)\) 。

\begin{enumerate}
\def\labelenumi{\arabic{enumi})}
\setcounter{enumi}{19}
\item
设 \((u_n)\) 是有限实数 \(\geqslant 0\) 的一个序列，满足 \(u_0 > 0\) ，并令 \(s_n = \sum_{k=0}^{n} u_k\) （对每个 \(n \geqslant 0\) ）。则序列 \((u_n)\) 在 \(\mathbf{R}\) 中可和，当且仅当序列 \((u_n / s_n)\) 在 \(\mathbf{R}\) 中可和{[}对序列 \((1 - u_n / s_n)]\) 应用定理 4。若把序列 \((u_n / s_n)\) 换成序列 \((u_n / s_{n-1})\) ，结论同样成立。
\item
令 \(u_{n} = (-\mathrm{I})^{n} / \sqrt{n}\) （对 \(n \geqslant 2\) ）。以 \(\mathrm{I} + u_{n}\) 为通因子的乘积不收敛，但通项为 \(u_{n}\) 的级数收敛。
\item
对 \(n \geqslant 2\) ，令 \(u_{2n-1} = -\mathrm{I}/\sqrt{n}\) 且 \(u_{2n} = (\mathrm{I} + \sqrt{n})/n\) 。以 \(\mathrm{I} + u_n\) 为通因子的乘积收敛，但通项为 \(u_n\) 的级数不收敛。
\end{enumerate}

